\begin{document}
\begin{titlepage}
\vspace*{13mm}
{\Large\color{teal} TFPT \quad / \quad Universalraum \quad / \quad Prozess und Physik\par}
\vspace{16mm}
{\Huge\bfseries\color{navy} Die gemeinsame Quelle\par}
\vspace{5mm}
{\LARGE und der Weg zur Gesamtlösung\par}
\vspace{13mm}
{\large Gesamtsynthese der Sitzung vom 28. September 2026\\Fortgesetzt bis zum 30. September 2026\par}
\vspace{8mm}
\key{\textbf{Das Ganze: ein Ursprung, mehrere gekoppelte Antworten.}\par
TFPT besitzt bereits konkrete Herleitungsketten für Quellenmassen, Yukawa-Hierarchien, Mischungen und $\alpha$ sowie bedingte Neutrino-, Gravitations- und Kosmologieanschlüsse. Der gemeinsame E8-Quellkandidat liefert dazu eine ausführbare Algebra mit Zustand, Feldprodukten und Konformalzeit. Die fundamentale Aufgabe ist, beide Seiten als Wirkungen desselben ursprünglichen Prozesses zu identifizieren.\par
Die erneute Prüfung unterscheidet die Auswahl einer gemeinsamen physischen Antwortregel von der Wahl ihrer mikroskopischen Darstellung. Der kurze Viersektor-/Fixpunktvorschlag vom 30. September wird auf den folgenden Seiten am Original und an exakter Arithmetik geprüft. Er komprimiert die Verklebung, ersetzt aber weder den tatsächlichen Rang acht noch Zustand, räumlichen Zugriff und Dynamik. Für die verallgemeinerte Zwei-Seiten-Struktur $A_{d-1}\oplus D_{d+1}$ bleibt außerdem der allgemeine bedingte Satz erhalten: Ein einziger lokaler Spinorstrom vom konformen Gewicht eins wählt ausschließlich $d=4$ und E8. Dieser Feldkandidat ist ein möglicher Anschluss, nicht das vorab bestimmte Zentrum jedes weiteren Lösungswegs.\par
\textbf{Die erneute Gesamtprüfung korrigiert einen konkreten Suchfehler:} Die vorhandene endliche Clock-Vergröberung ist eine Auslese, keine produktbewahrende Darstellung der gesamten Quelle. Unsichtbare Einzelschritte können zusammen eine sichtbare Antwort erzeugen. Die vollständige graduierte Operationskette muss deshalb vor der Auslese erhalten bleiben. Der ausgewählte QWZ-Kragen liefert bereits bedingte kontinuierliche neutrale Antworten und geladene CAR-Felder mit Vakuum und Dynamik. Seine Auswahl aus P1/P2 und der gemeinsame physische Anschluss sind damit noch nicht bewiesen. Die bestehenden Massen-, Flavor- und Alpha-Herleitungen bleiben erhalten; gesucht wird die ursprüngliche gemeinsame Operations-, Zustands- und Anschlussregel, die ihre Antworten zusammen erzeugt.}
\vspace{8mm}
Dieses Dokument führt die ursprünglichen TFPT-Paper, die Konsolidierungen vom 26.--28. September, die eingefügten Synthesevorschläge und die ausführbaren Ergebnisse dieser Sitzung zusammen. Es erklärt die Kette in einfacher Sprache und mit Skizzen, erhält die mathematischen Voraussetzungen und entwickelt daraus konkrete Forschungs- und Algorithmuswege.

\vfill
\textbf{Für Stefan Hamann}\par
Arbeitsfassung mit Quellen- und Befundregister. Die geometrischen Zeichnungen sind erklärende Schemata; sie behaupten keine bereits abgeleitete Gestalt der Raumzeit. Der Anwendungsteil interpretiert das gesprochene Wort ``hallieren'' vorläufig als Hylæan/Field Intelligence.
\end{titlepage}

\section*{Korrektur der Suchrichtung am 29. September}
\addcontentsline{toc}{section}{Korrektur der Suchrichtung am 29. September}
\key{Der Nutzer hat zu Recht beanstandet, dass die Untersuchung immer wieder dieselbe offene Wirkungsfrage reformuliert. Der zuletzt konstruierte E8-Quellkandidat wird deshalb als vorab gewähltes Zentrum der weiteren Suche zurückgenommen. Seine mathematischen Ergebnisse bleiben erhalten; ihre Existenz bestimmt nicht die Forschungspriorität.}

Der Originalaufbau beginnt mit der verbundenen Naht-/Trägerstruktur. In
\file{introduction.tex:140-290} und \file{introduction.tex:355-433} sind diskreter Abschluss und kontinuierliche Randantwort zwei verbundene Wege; E8 ist ihr algebraischer Abschluss und Konsistenzrahmen. Der E8-Randnetzansatz gehört selbst zum Original, doch seine Gleichsetzung mit der rohen Naht ist Gegenstand des bedingten Vertrags \file{SEAM.EQUIV.01}. Einen ausgewählten ausführbaren Randnetzkandidaten auszubauen beweist diese Gleichsetzung nicht und ist keine vorab erzwungene Fortsetzung des Gesamtauftrags.

Eine weitere Korrektur betrifft Abschnitt~\ref{sec:commonvariation}: Der Vertrag
\file{FTRANSFER.GENERATING.01} verlangt einen gemeinsamen vierdimensionalen Erzeuger für vier nachgelagerte Transfers (Polmassen, Baryonasymmetrie, Axionrelikt und Proton-/Elektronmassenverhältnis), sofern die entsprechende Theorie mit kontrolliertem Kontinuum vorliegt. Er identifiziert weder diesen Erzeuger mit der dort betrachteten endlichen Operatorfamilie noch setzt er die bereits vorhandenen Compiler-Herleitungen außer Kraft. Die vorherige Übertragung seiner Aussage auf den gesamten Ursprung war zu weitgehend.

Der erneute Ausgangspunkt ist daher der \emph{gesamte ursprüngliche Zusammenhang mit seinen Rückbedingungen}: Naht und Doppeldeckung, vier Marken, Träger und Familien, Clocks, Flavor, Alpha sowie der geometrische Randkanal. Die neueren Code-, Quellen- und Graphmodelle sind daran zu prüfende Kandidaten. Ein gemeinsamer Ursprung kann verschiedene Räume und wirksame Beschreibungen besitzen; er verlangt nachgewiesene verträgliche Abbildungen, nicht ohne Begründung dieselbe endliche Matrix für alle Erscheinungen. Die verbleibenden Beweisaufgaben sind damit nicht gelöst. Zurückgenommen wird die unbegründete Behauptung, der zuletzt verfolgte Operatoranschluss sei bereits als stärkster oder einzig richtiger Gesamtweg ausgewählt.

\section*{Kurswechsel am 30. September: Der einfache Kern und seine vollständige Antwort}
\label{sec:fixedpointproposalcheck}
\addcontentsline{toc}{section}{Kurswechsel am 30. September: Der einfache Kern und seine vollständige Antwort}

\key{\textbf{Die sinnvollere Frage lautet jetzt: Welche kleinste gemeinsame Regel erzeugt alle ursprünglichen Antworten?} Die bisherige Suche nach einem bestimmten mikroskopischen Operator darf nicht voraussetzen, dass dessen Matrixdarstellung physisch eindeutig sein muss. Eine vollständige Antwortregel kann Quelle, Operatoren und Zustand bis auf eine gemeinsame unitäre Umbenennung bestimmen. Endliche Invarianten allein leisten dies nicht. Die Vermutung, TFPT beschreibe eine universelle Fixpunktklasse vieler mikroskopischer Realisierungen, ist eine ernstzunehmende alternative Interpretation; ihre Universalität ist noch nachzuweisen.}

\subsection*{Was der neue Viersektorvorschlag tatsächlich beweist}
Für die ungeraden $D_g$-Wurzelgitter besitzt die Diskriminantengruppe die Form $\Z_4$. In der Konvention für konforme Spins, also Normquadrat geteilt durch zwei, gilt
\[
 q_g(a)=\frac{ga^2}{8}\pmod1,\qquad
 \frac12\sum_{a=0}^3e^{2\pi iq_g(a)}=e^{2\pi ig/8}.
\]
Die Form auf $A_3$ lautet $q_{A_3}(a)=3a^2/8\pmod1$. Die diagonale Glue-Untergruppe ist genau dann isotrop, wenn
\[
 \boxed{g+3\equiv0\pmod8.}
\]
Dies ist die kurze ursprüngliche Diskriminantenverklebung, jetzt in der Spin-Konvention ausgeschrieben; Abschnitt~\ref{sec:globaloriginselection} verwendet die doppelte Norm-Konvention modulo zwei. Isotropie plus die tatsächliche Index-vier-Verklebung liefert in dieser Klasse einen geraden unimodularen Abschluss~\cite{nikulinforms}.

\begin{center}
\begin{tabular}{rllr}
\toprule
$g$ & $q_g(0),q_g(1),q_g(2),q_g(3)$ & Diagonal verklebbar & Rang $g+3$\\\midrule
3 & $0,3/8,1/2,3/8$ & nein & 6\\
5 & $0,5/8,1/2,5/8$ & ja & 8\\
7 & $0,7/8,1/2,7/8$ & nein & 10\\
13 & $0,5/8,1/2,5/8$ & ja & 16\\\bottomrule
\end{tabular}
\end{center}

Die Tabelle wurde erneut symbolisch gerechnet. Sie macht zugleich die Grenze sichtbar: $g=5$ und $g=13$ besitzen \emph{dieselbe} endliche quadratische Form. Die Form wählt $g\equiv5\pmod8$, nicht schon $g=5$. Der zusätzliche \emph{tatsächliche} Rang acht wählt $g=5$ und E8. In der vorausgesetzten Gitter-CFT ist $c$ gleich dem Rang; $c\bmod8$ enthält dagegen gerade diese Information nicht. Für stark rationale holomorphe Vertexoperatoralgebren mit tatsächlichem $c=8$ gibt es außerdem den einschlägigen Klassifikationssatz von Dong--Mason~\cite{dongmasonlow}. Dessen Lokalitäts-, Rationalitäts- und Holomorphievoraussetzungen dürfen nicht durch die endliche Form ersetzt werden.

Die vier Diskriminantenklassen sind zunächst \emph{Ladungssektoren}, nicht vier Punkte eines räumlichen Randes. Die Abbildung auf die geometrischen Marken ist Teil des ursprünglichen TFPT-Wörterbuchs. Nach der unimodularen E8-Erweiterung ist die volle Diskriminantengruppe trivial; die vier Klassen bleiben als Zerlegung bei Einschränkung auf $D_5\oplus A_3$ erhalten. Ihre gemeinsamen Spins sind ganzzahlig: $5/8+3/8=1$ und $1/2+1/2=1$. Daher liefert $q$ allein keine vierdimensionale Spin-Ladungs-Relation. Insbesondere sind die geklebten Gewicht-eins-Ströme bosonisch, obwohl einer ihrer internen Faktoren ein Spinor ist; siehe Abschnitt~\ref{sec:physicaltimeaudit} und den Materieanschluss.

\subsection*{Die vorgeschlagene Ordnung 24 und die Ratenprüfung}
Für die gewöhnlichen affinen Charaktere ist die modulare Matrix diagonal,
\[
 T_{aa}=e^{2\pi i(h_a-c/24)}.
\]
Die folgenden Ordnungen wurden aus diesen rationalen Exponenten exakt berechnet. Die relative Matrix ist $T/T_{00}$.
\begin{center}
\begin{tabular}{lrrl}
\toprule
Theorie & Ordnung $T$ & Ordnung $T/T_{00}$ & Grundgewichte\\\midrule
$D_{5,1}$ & 24 & 8 & $0,5/8,1/2,5/8$\\
$A_{3,1}$ & 8 & 8 & $0,3/8,1/2,3/8$\\
$D_{5,1}\times A_{3,1}$, vor Glue & 24 & 8 & sechzehn Produktsektoren\\
Geklebtes $E_{8,1}$ & 3 & 1 & nur der Vakuummodul\\\bottomrule
\end{tabular}
\end{center}
Auf den diagonalen Glue-Sektoren hat das Produkt denselben skalaren Faktor $e^{-2\pi i/3}$. Die $24$ stammt in der $D_5$-Charakterkonvention auch aus der Vakuumphase $-5/24$; sie ist keine neu hergeleitete 24-Schritt-Bewegung des gesamten E8-Systems. Die gespeicherte andere Seam-Clock $T_{\rm seam}=R\otimes G_F$ aus \file{universalraum-seam-closure-20260915/RESULTS.md:33-52} hat ebenfalls absolute Ordnung 24, aber $T_{\rm seam}^{12}=-I$ und kleinste projektive Periode 12. Dagegen ist $T_{D_5}^{12}$ in der Sektorreihenfolge $0,s,v,c$ gleich $\diag(-1,1,-1,1)$ und nicht skalar; seine kleinste projektive Periode ist acht. Diese beiden Clocks können also trotz der gleichen absoluten Ordnung nicht durch einen Basiswechsel und eine gemeinsame skalare Phase identifiziert werden. Der ursprüngliche Clock-Vertrag selbst behauptet keine Gleichsetzung mit einer Beobachtungsuhr.

Ein Dehn-Twist wirkt auf konformen Blöcken beziehungsweise auf der modularen Beschreibung des Torus. Seine Identifikation mit Laborzeit benötigt eine eigene Abbildung. Der konforme Generator $L_0$ liefert tatsächlich Zylinderzeit, aber keine automatische Gleichheit von RG-Skala, modularer Transformation und Minkowski-Translation. Abschnitt~\ref{sec:physicaltimeaudit} schließt die exakte affine Gleichsetzung einer periodischen $L_0$-Rotation mit nichttrivialer globaler Poincar\'e-Zeit bereits aus; lokale Grenzübergänge und andere Raumzeitlesarten bleiben davon getrennt.

Auch der konkrete Test der Zahlen $8/5,9/64,10/9,41/10$ ist entschieden: Die Grundgewichte von $D_{5,1}$ sind $0,1/2,5/8$, die von $A_{3,1}$ sind $0,1/2,3/8$. Nachfahren addieren ganze Zahlen. In der Produktquelle besitzen daher alle Gewichte Nenner, die acht teilen; lokale E8-Gewichte sind ganzzahlig. \emph{Keine} der vier genannten Zahlen ist ein solches konformes Gewicht. Das widerlegt diese direkte Spektrumsidentifikation, nicht mögliche deformierte Theorien, andere affine Level oder Renormierungsantworten. Ein Betakoeffizient wie $41/10$ und eine Relaxationsrate sind schon ihrem Typ nach nicht automatisch ein Feldgewicht.

Ihre tatsächlichen Rollen sind verschieden: $10/9$ ist das Produkt der Verzweigungspunkte $-2/3$ und $-5/3$ des Flavor-Pencils, \file{tfpt_2_standard_model.tex:1397-1420}. $41/10$ ist der hyperladungsnormalisierte Ein-Schleifenkoeffizient aus dem konkreten Drei-Familien-plus-Higgs-Inhalt, \file{tfpt_2_standard_model.tex:1638-1649}. $9/64$ tritt in \file{v508_celestial_wp5e_gamma_pairing.py:51-56,576-596} als Kopplung eines bereits ausgeschlossenen Sphere-Axion-Austauschansatzes auf. Für $8/5$ wurde in der gezielten Suche kein eigenständiger TFPT-/SM-Ratensatz gefunden; ein exakter Fund liegt im gesonderten Prime/RH-Strang. Die vier Zahlen bilden damit keine bereits belegte gemeinsame Liste von dynamischen TFPT-Raten. Die gezielte Suche behauptet keine vollständige semantische Auswertung aller Dateien.

\subsection*{Was der funktionierende Rest über den Ursprung aussagt}
Die Trennung ``alles Erfolgreiche ist invariant, alles Dynamische scheitert'' ist am Original zu grob. Vorhanden sind exakte Transportiteration, positive Markov-Ausführung, unitäre Wortoperationen, KMS-Rekonstruktion sowie in dieser Sitzung analytische lokale Quellgrenzwerte. Ihre \emph{gemeinsame Herkunft und physische Identifikation} bleibt zu zeigen. Das ist eine andere Aufgabe als die Existenz irgendeiner Dynamik.

Die sektorübergreifenden Formeln sind dabei wertvolle Rückbedingungen. Beispielsweise bestimmt ein positiver Seed $\phi_0$ in der vorhandenen Klasse $\phi(c)=4c/3+48c^4$ wegen $\phi'(c)>0$ genau ein $c$; die bekannten Flavor- und Alpha-Lesarten können nicht unabhängig nachjustiert werden. Gleichzeitig enthalten Tausende Prüfungen viele Folgen derselben Eingaben. Ihre Anzahl zählt deshalb nicht die Zahl unabhängiger Bedingungen an den gesamten Quelloperator. Die bestehenden Massen-, Yukawa- und Neutrinoherleitungen werden durch diese Feststellung nicht zurückgenommen.

Die älteren Gegenfälle zeigen konkret, \emph{welche} Information häufig früh verloren wurde. Ein Signprojektor behält das Vorzeichen eines Generators, nicht seine Energieabstände; \file{seam-state-dynamics-injectivity-20260918/README.md} beweist dies. Klassische Transferdaten behalten Populationen, nicht sämtliche Phasen oder das Gedächtnis der Umgebung. Der Originaltext in \file{origin_theory.tex:780-796} unterscheidet selbst mehrere Quantisierungen desselben klassischen Transports. Die bloße Forderung nach einer zusätzlichen möglichst einfachen Matrix repariert diesen Informationsverlust nicht.

Der alte operative Ursprung enthält diese reicheren Daten bereits als \emph{Eingaben}:
\[
 S_{\rm op}=(\mathcal A_{\rm loc},\tau_t,\Theta,\omega,[u_\Sigma],D_{\rm coll}).
\]
\file{_archive/tfpt-45/source_extracts/01_boundary_kernel_source.tex:480-513} setzt unter anderem die Zeitaktion $\tau_t$ und den reflexionspositiven Zustand $\omega$ voraus und bildet daraus das Randdatum. Das ist der konkret übersehene Unterschied: Die alte Schnittstelle ist dynamisch, aber sie \emph{leitet} diese Eingaben nicht schon aus dem endlichen Compiler ab. Ebenso ist nicht bewiesen, dass ihre anschließende Kompression zu Rand-/Fixpunktdaten Zeit und Zustand vollständig rückgewinnbar erhält. Die Antworttreue dieser ursprünglichen Abbildung ist deshalb ein präziserer Forschungsgegenstand als die Suche nach einer weiteren freien Matrix.

Ein vorhandener Zeuge wurde in dieser Prüfung erneut ausgeführt: \file{compiler-origin-audit-20260913/label_lift.py:160-176} vergleicht zwei Realisierungen derselben klassischen $15$-Label-Randmatrix und derselben sequentiellen Labelgeschichten. Bei kohärentem Eingang und derselben Projektion beträgt die Antwort $1/15$ beziehungsweise $1$. Die damalige Operations-/Auslesewahl ist dabei Voraussetzung, nicht neu hergeleitete TFPT-Physik. Der Zeuge zeigt genau, dass dieser gemeinsame klassische Schatten nicht die vollständige Antwort bestimmt. \file{universalraum-inversion-20260914/checker.py:172-203} ergänzt einen gespeicherten Fortsetzungszeugen: Gleiche momentane reduzierte Zustände liefern bei Wiederverwendung desselben gekoppelten Registers unterschiedliche Antworten. Die alte Schnittstelle darf deshalb nicht mit einem gedächtnislosen Gesamtprozess identifiziert werden.

\subsection*{Der andere Lösungsweg: Die Antworttreue zum Ausgangspunkt machen}
\picturetext{Mehrere Programme können dieselbe Liste von Ergebnissen ausgeben. Ob sie dieselbe Regel ausführen, erkennt man daran, was nach Eingriffen passiert: dieselben Schritte anders anordnen, eine Phase vergleichen oder einen früheren Speicher nochmals benutzen. Der Ursprung wird durch diese Fortsetzungen bestimmt, nicht durch seine Dateiform oder die Namen seiner Matrizen.}

Gesucht wird deshalb die kleinste \emph{begründete} Operationsregel mit vollständigen Antworten. Für eine festgelegte $*$-Wortalgebra, einen positiven normierten Zustand und einen zyklischen, von diesen Worten erzeugten Raum gilt
\[
 K(u,v)=\omega(u^*v).
\]
Der Quotient der Worte nach den Nullvektoren von $K$ bestimmt Hilbertraum, Vakuum und die markierten Operationen bis auf eine gemeinsame unitäre Identifikation, sofern Linksverkettung wohldefiniert ist und die erforderlichen Operator-/Domänenbedingungen erfüllt sind. Eine beidseitige Zeitgruppe von $*$-Automorphismen, die den Zustand und damit diesen Kern erhält, rekonstruiert zusätzlich ihre unitäre Wirkung; bei starker Stetigkeit folgt ihr selbstadjungierter Generator. Eine bloße kernerhaltende einseitige Verschiebung liefert zunächst nur eine Isometrie. Dieser Rückweg wurde in den früheren Wort-/KMS-Verträgen bereits bewiesen. Er ist keine neue Herkunftsherleitung: Das physisch zulässige Alphabet, seine Lokalisation und die Antwortregel selbst müssen aus den Originalprinzipien folgen.

Die nächste Untersuchung soll deshalb nicht wieder eine beliebige Quelle wählen und anschließend alle Ausgaben daran anpassen. Sie prüft die \emph{Faser des ursprünglichen Compilers}: Welche vollständigen Prozesse erfüllen gleichzeitig seine markierten Beziehungen, seine ursprünglichen Auslesungen und die schon vorhandenen physikalischen Herleitungen? Unterschiede, die durch einen gemeinsamen Basiswechsel verschwinden, sind keine physikalisch verschiedenen Ursprünge. Unterschiede in einer zulässigen gemeinsamen Antwort müssen dagegen durch die Originaldaten entschieden werden. Lokale Feldumdefinitionen können ebenfalls physikalische Observablen erhalten; die stärkere Wort-/Quellenäquivalenz erfordert jedoch den passenden Transport aller Einsetzungen und darf nicht aus bloßer S-Matrix-Äquivalenz gefolgert werden~\cite{criadoredefinitions}.

Eine zielführende Ablation variiert deshalb zuerst vollständige Kandidaten \emph{bei festgehaltenen Invarianten}, statt die Zutaten einer Formel nacheinander zu löschen. Die Prüfungsreihe lautet: gleiche endliche Quadratform und gleiche strukturelle Compilerantworten; dann dieselbe kohärente Fortsetzungsfrage auf beiden Seiten; anschließend derselbe gemeinsame Flavor-/Alpha-/Skalenanschluss einschließlich gemischter Variationen. Werden verschiedene Antworten gefunden, ist genau die trennende Originaloperation zu identifizieren. Bleiben alle begründeten Antworten gleich, ist eine noch verbleibende Darstellungswahl kein Grund, einen besonderen Mikrounterbau zu erzwingen. Die Behauptung einer gemeinsamen Universalitätsklasse verlangt darüber hinaus einen kontrollierten Vergröberungs- beziehungsweise Grenzsatz.

Für den Vergleich zweier \emph{vollständiger} Quellenrepräsentationen ist eine explizite Antworttreue-Bedingung sinnvoll. Für die tatsächlichen quellseitigen Operationen und eine nicht nachträglich passend gewählte volle Darstellung $\Phi$ müssen gemeinsam gelten
\[
\begin{aligned}
\Phi(xy)&=\Phi(x)\Phi(y),&
\Phi(\Theta x)&=\Theta_{\rm gr}\Phi(x),\\
\Phi(\tau_t x)&=\alpha_t\Phi(x),&
\omega_{\rm gr}(\Phi x)&=\omega(x).
\end{aligned}
\]
Graduierung, Windung und Lokalisation bleiben dabei markiert; zusätzlich ist $\Phi(x^*)=\Phi(x)^*$ erforderlich. \textbf{Diese Bedingungen gelten nicht für die vorhandene endliche Clock-Vergröberung $E_0$.} Die bisherige Formulierung hatte eine volle Darstellung und eine Auslese zu wenig unterschieden. $E_0$ ist eine bedingte Erwartung und im Allgemeinen nicht multiplikativ. Abschnitt~\ref{sec:sourcedirectionaudit} korrigiert diese Suchbedingung am Original. Ein geordnetes Vierpunktwort ist weiterhin ein sinnvoller erster Trennversuch. Ein bestandener Einzeltest beweist jedoch weder sämtliche Wortantworten noch die vollständige Raumzeitquelle; dafür wird der vollständige Kern beziehungsweise ein begründeter Erzeugungs- und Relationensatz benötigt.

Der physikalische Teil beginnt mit einer notwendigen gemeinsamen Anschlussbedingung: Auf \emph{ursprünglich identifizierten} überlappenden Observable-/Feldräumen müssen die verschiedenen Sektorantworten dieselbe Kovarianz und dieselbe Quellenvariation liefern. Schematisch müssen bei feststehenden Einbettungen $V_i$ die Daten $C_i$ Kompressionen \emph{einer} positiven $C$ sein,
\[
 C_i=V_i^*CV_i.
\]
Einbettungen dürfen nicht nachträglich disjunkt gewählt werden, um jeden Widerspruch zu verstecken. Außerdem ist im Allgemeinen $\log(V_i^*CV_i)\ne V_i^*(\log C)V_i$; Kompression, effektive Zeit und Gedächtnis sind gemeinsam zu erhalten. Der inverse Kompatibilitätstest kann Kandidaten ausschließen, beweist aber noch nicht die Herkunft der vervollständigenden Kovarianz. Der gemeinsame Determinantenanschluss mit Verbindung ist bereits \file{SEAM.DETLINE.UNIFICATION.01}; \file{FTRANSFER.GENERATING.01} betrifft enger die vier nachgelagerten Transfers, nicht automatisch sämtliche Compilersektoren. Die benötigten gemischten Variationen stehen in Abschnitt~\ref{sec:commonvariation}. Kein neuer universeller Einheitsclock oder beliebiger Hamiltonoperator wird dafür vorausgesetzt.

\key{\textbf{Arbeitsurteil:} Der Viersektorvorschlag komprimiert einen echten algebraischen Kern. Seine Gleichsetzung mit der vollständigen Quelle, die Deutung der $24$ als physischer Takt und die vier vorgeschlagenen Gewichtstreffer tragen nicht. Der fundamentale Kurswechsel besteht darin, das vorhandene Gesamtgefüge über seine \emph{gemeinsamen Operationen und Antworten} zu bestimmen. Die Einfachheit einer Grundregel ist mit einer anspruchsvollen Grenzherleitung vereinbar; aus der bisherigen Schwierigkeit folgt weder, dass der Ursprung kompliziert sein muss, noch, dass die fehlende Antwort bereits durch die endlichen Invarianten festgelegt ist.}

\section*{Erneute Gesamtprüfung: Die Auslese ist nicht die ganze Maschine}
\label{sec:sourcedirectionaudit}
\addcontentsline{toc}{section}{Erneute Gesamtprüfung: Die Auslese ist nicht die ganze Maschine}

\key{\textbf{Der konkrete Kursfehler ist nachgewiesen.} Die bisherige Suche verlangte stellenweise, dass eine neutrale Vergröberung alle Produkte der vollständigen Quelle wie eine Darstellung bewahrt. Die vorhandene endliche Clock-Vergröberung tut dies ausdrücklich nicht. Schon die vorhandene Prüfung \file{v1030} enthält den passenden Gegenfall. Der nächste Forschungsgegenstand ist daher die ursprüngliche \emph{vollständige graduierte Operationskette vor ihrer Auslese}, zusammen mit Zustand, Orientierung, Zeit und dem gemeinsamen physikalischen Anschluss. Diese Korrektur liefert keine vollständige physikalische Ursprungslösung; sie entfernt eine falsche Bedingung aus der Suche.}

\subsection*{Was über den gesamten Bestand geprüft wurde}
Die aktuelle Bestandsprüfung erfasst die Register und Einordnungen von 1028 Verifikationsskripten, 279 Theorie-Contracts, 59 Experimentverzeichnissen und 179 Lean-Dateiknoten einschließlich RH. Das ist keine Behauptung, jedes Skript neu ausgeführt oder jede Zeile vollständig semantisch geprüft zu haben. Die gezielte Originalprüfung verfolgt die tragenden Wege: P1/P2 und Compiler, Clock und Vergröberung, Quellen-/Determinantenanschluss, Flavor und Neutrinos, Universalraum sowie externe und interne Experimente. Besonders entscheidende Gegenfälle und aktuelle Rückbedingungen wurden erneut ausgeführt. Die Herkunftsaussagen unten beziehen sich auf die benannten Originale, nicht auf die Größe des Registers.

Die 127 Zeilen der Evidenz-Scorecard unterscheiden unter anderem 32 Vergleiche mit externen Daten, 20 nachgelagerte physikalische Übertragungen, 60 Suchziele und 12 interne Konsistenzprüfungen. Ein interner KMS-Test, eine erfolgreiche Instrumentenkalibrierung und eine Übereinstimmung mit einer gemessenen Größe entscheiden unterschiedliche Fragen. Sie dürfen weder als gleichartige Bestätigungen addiert noch pauschal als bedeutungslos verworfen werden.

Auch der Theoriegraph ist kein automatischer Beweisgraph. Der ursprüngliche Vertrag \file{GATE.LEDGER.01} bezeichnet seine Abhängigkeitspfeile als gegenseitige Evidenz- und Querverweise. Eine geschlossene Schleife solcher Pfeile beweist keine gegenseitige Herkunft. Entscheidend bleiben die Voraussetzungen und das Verdict des jeweiligen Originals.

\subsection*{Der kleinste entscheidende Test an der vorhandenen Clock-Vergröberung}
In \file{verification/v993_minimal_defect_selector.py:169-212} ist die Clock
\[
 u=\diag(1,i,-1,-i,1,i,-1,-i,\ldots)
\]
und ihre neutrale Auslese
\[
 E_0(x)=\frac14\sum_{k=0}^{3}u^k x u^{-k}.
\]
Die Originalprüfung bestätigt Idempotenz, Positivität, Spurerhaltung und die Regel
\[
 E_0(b_1xb_2)=b_1E_0(x)b_2,\qquad b_1,b_2\in\mathcal A_0.
\]
Das ist die Struktur einer bedingten Erwartung. Die Faktoren außerhalb von $x$ müssen dabei bereits neutral sein. Eine allgemeine Produktregel $E_0(xy)=E_0(x)E_0(y)$ folgt daraus nicht~\cite{tomiyamaexpectation}.

Der in dieser Prüfung frisch ausgeführte Originalzeuge nimmt im 16-dimensionalen Clockraum
\[
 X=\ket{1}\bra{0},\qquad X^*=\ket{0}\bra{1},\qquad P_0=\ket{0}\bra{0}.
\]
Die Indizes bezeichnen hier \emph{interne Clocksektoren}; sie sind keine bereits hergeleiteten Orte oder Teilchen. Exakt gilt
\[
 E_0(X)=E_0(X^*)=0,\qquad
 E_0(X^*X)=P_0,\qquad E_0(X^*)E_0(X)=0.
\]
Der Unterschied hat Operatornorm eins; seine normierte Spur beträgt $1/16$. Das zusammengesetzte Wort $X^*X$ bleibt sichtbar, obwohl beide einzelnen Faktoren in der neutralen Auslese verschwinden.

\begin{center}
\begin{tikzpicture}
\node[box,text width=36mm] (full) {Beide vollständigen\\Schritte behalten\\$X,\ X^*$};
\node[box,text width=36mm,right=5mm of full] (word) {Zuerst zusammensetzen\\$X^*X=P_0$};
\node[box,text width=36mm,right=5mm of word] (read) {Dann neutral auslesen\\$E_0(P_0)=P_0$};
\node[box,text width=36mm,below=7mm of full] (early) {Beide einzeln auslesen\\$E_0(X)=E_0(X^*)=0$};
\node[box,text width=36mm,right=5mm of early] (zero) {Auslesungen multiplizieren\\$0\cdot0=0$};
\node[box,text width=36mm,right=5mm of zero] (lost) {Die gemeinsame\\Antwort ist verloren};
\draw[arr] (full)--(word);\draw[arr] (word)--(read);
\draw[arr] (early)--(zero);\draw[arr] (zero)--(lost);
\end{tikzpicture}
\end{center}

\picturetext{Ein Übergang verlässt die sichtbare Ebene und ein zweiter kehrt zurück. Wer jeden Schritt sofort wegfiltert, kann die Rückkehr nicht mehr berechnen. Das Bild beschreibt die Algebra des Originals; es ist keine zusätzliche Behauptung über räumliche Bewegung.}

Der ursprüngliche \file{v1030_alg2_sector_frame.py:74-106} warnt mit einem ähnlichen Projektionszeugen ausdrücklich davor, tatsächlich ausgeführte Worte durch Produkte ihrer Kompressionen zu ersetzen. Die jetzige Korrektur entdeckt daher keine fehlende Eigenschaft des Compilers. Sie macht eine bereits vorhandene Warnung wieder zum Maßstab der Suche. Allgemein zeigt die Bimodulregel
\[
 E_0(xy)-E_0(x)E_0(y)
 =E_0\bigl((x-E_0(x))(y-E_0(y))\bigr).
\]
Die wegprojizierten Anteile können gemeinsame neutrale Wirkungen haben. Der Kern von $E_0$ ist deshalb im Allgemeinen kein Ideal, nach dem man sämtliche Operationen folgenlos quotientieren könnte.

\subsection*{Welche einfache vollständige Struktur tatsächlich erhalten werden muss}
Für das vorhandene Clock-Automorphismusdatum gibt es die vier Gradprojektionen
\[
 E_r(x)=\frac14\sum_{k=0}^{3}i^{-rk}u^kxu^{-k},\qquad
 x=\sum_{r\in\Z_4}E_r(x).
\]
Die ganze Zerlegung behält alle Operationen. Nur $E_0$ ist dabei eine positive bedingte Erwartung; die übrigen $E_r$ sind Spektralprojektionen, keine positiven Auslesungen. \emph{Keine einzelne Projektion wird deshalb multiplikativ.} Die korrekte Produktregel lautet
\[
 E_t(xy)=\sum_{r\in\Z_4}E_r(x)E_{t-r}(y),\qquad
 E_r(x^*)=E_{-r}(x)^*.
\]
Im Zeugen ist $X$ vom Grad eins, $X^*$ vom Grad drei und ihr Produkt neutral. Der operatorwertige Kern $K_E(x,y)=E_0(x^*y)$ behält ihre gemeinsame Wirkung. Er wird aus der vollständigen Abbildung mit ihrem Definitionsbereich gebildet; die neutrale Bildalgebra allein bestimmt ihn nicht.

Der richtige Vergleich hat somit zwei verschiedene Pfeile:
\[
 \mathcal A_{\rm src}\xrightarrow{\ \Phi\ }\mathcal A_{\rm gr}
 \xrightarrow{\ E_0\ }\mathcal A_0.
\]
Die volle Darstellung $\Phi$ muss tatsächliche Produkte, Adjungierte, Markierungen, Zeit und Lokalisation verträglich übertragen. $E_0$ wird anschließend auf das \emph{ganze Wort} angewendet. Für einen neutralen Gesamtgrad wirkt $E_0$ identisch; die Zwischenfaktoren können dennoch geladen sein. Eine Auslese nach jedem Zwischenfaktor beschreibt eine andere Operation.

Zustandstreue ist eine weitere echte Bedingung:
\[
 \omega_{\rm src}(w)=\omega_{\rm gr}(\Phi(w))
 =\omega_0\bigl(E_0(\Phi(w))\bigr).
\]
Die letzte Gleichheit setzt $\omega_{\rm gr}=\omega_0\circ E_0$ voraus. Sie gilt bei bewiesener Invarianz des vollen Zustands unter der Deckwirkung, mit $\omega_0=\omega_{\rm gr}|_{\mathcal A_0}$. Das ist eine Faktorisierung durch die Erwartung, keine Produktzustandsannahme. Sie folgt nicht schon aus Positivität der Vergröberung und muss aus dem ursprünglichen Zustand bewiesen werden. Für die Zeitverträglichkeit muss die Zeitwirkung zusätzlich mit der Deckwirkung kommutieren, damit $E_0\alpha_t=\alpha_tE_0$. Die ursprüngliche volle Antwortregel darf daher nicht durch beliebig gewählte neutrale Zustandsdaten ersetzt werden.

Diese Unterscheidung passt auch zum Universalraum: Eine momentane ausgelesene Ansicht ist nicht automatisch der gesamte fortsetzbare Prozess. Der ursprüngliche Zeuge in \file{universalraum-inversion-20260914} zeigt, dass eine erneut benutzte Umgebung Antworten unterscheiden kann, obwohl der momentane reduzierte Zustand gleich ist. Vollständige Mehrzeitdaten können solche Prozesse charakterisieren; dies ist auch der Gegenstand des Prozess-Tensor-Formalismus~\cite{pollockprocess}. Daraus folgt weder automatisch eine TFPT-Ursprungsregel noch eine festgelegte Raumzeit.

\subsection*{Was die tragenden Verifikationen wirklich schließen}
\begin{tabularx}{\linewidth}{L{29mm}Y}
\toprule
Original & Belastbare Reichweite und verbleibender Anschluss\\\midrule
\file{v991}, \file{v1014} &
Bulk-/Randindex, Orientierung und endliche Determinantenlinie werden gemeinsam abgeglichen. Eine konstante Referenzphase und der kontinuierliche Anschluss mit Verbindung sind gesondert zu behandeln.\\
\file{v1015} &
Im ausgewählten QWZ-Kragen liefern die ursprüngliche demokratische Antwort und ihre Mutanten einen echten Auswahltest. Der frische Lauf besteht 11 Prüfungen. Der Kragen ist damit nicht aus P1/P2 selbst ausgewählt; AX.P2 bleibt Ursprungsannahme.\\
\file{v1026} &
Der TEL-B-Normschritt ist für die benannte Familie gleichmäßig geschlossen. Die Originalgrenze lautet $2.995906<3$ für gerade $N\ge16$ bei festem $M=1$, $N_y=8$. Stromerzeugung und Algebraerschöpfung folgen nicht allein aus dieser Norm.\\
\file{v1030} &
Der frische Lauf besteht 27 Prüfungen. Volle neutrale Blöcke und ein zusammenhängender geladener Sektorgraph erzeugen im endlichen Modell die volle Algebra. Der gleichmäßige Grenzschluss braucht tatsächlich ausgeführte Worte, einen nicht kollabierenden gemeinsamen Rahmen und kontrollierte Ausgabetails für Worte und Adjungierte. Diese mikroskopischen Voraussetzungen werden ausdrücklich nicht als bewiesen ausgegeben.\\
\file{SeamResidualAxiom.lean} &
Die Datei wurde frisch erfolgreich geprüft. Das Abschlusslemma benutzt ausdrücklich \file{CollarHasInvariants}, \file{SeamIsE8}, \file{collar_invariants} und \file{crossedproduct_route_theorem}. Lean prüft die Schlusskette unter diesen Voraussetzungen; es leitet die Quellenidentifikation nicht durch Umwandlung des Beweises in ausführbaren Code her.\\
\bottomrule
\end{tabularx}

Die vorhandenen geladenen Register- und Ereignis-Clock-Verträge bestätigen bereits, dass vollständige Produkte und Fortsetzungen wesentlich sind. Insbesondere können zwei Operatorpaare dieselben charakteristischen Polynome sämtlicher $A+tB$ besitzen und trotzdem verschiedene gemischte Produkte liefern: \file{source-flavor-joint-dictionary-20260921} erhält im ursprünglichen Paar $P_1(A_0)B_0P_3(A_0)=-3E_{31}$; im vereinfachten Paar verschwindet dieses Produkt. Ein Spektrum allein ist deshalb kein vollständiges Flavor-Wörterbuch.

Eine unabhängige Schranke bleibt erhalten: \file{source-graded-gns-matter-gate-20260922} zeigt, dass das Zentrum von Spin(10) auf $\operatorname{End}(\C^{32})$ durch Konjugation trivial wirkt, auf dem benötigten nativen Materiesektor jedoch nicht. Eine äquivariante Materieabbildung aus diesem Ausgangsraum muss dann verschwinden. Das Wiederaufnehmen aller Clockgrade repariert diesen falschen Darstellungsraum nicht automatisch. Es restauriert auch nicht automatisch allgemeines zeitliches oder umweltgetragenes Gedächtnis. Interne Graduierung, physische Fermionparität und Eichladung dürfen weiterhin nur über begründete Abbildungen identifiziert werden.

\subsection*{Was die Experimente als Rückbedingungen hinzufügen}
Die vorhandene \file{seed-consistency}-Auswertung ist besonders relevant, weil sie mehrere Sektoren an einen gemeinsamen Seed bindet. Die Originalanalysen v4--v6 wurden in dieser Prüfung direkt mit ihren eingefrorenen Messdaten erneut ausgeführt:
\begin{itemize}
\item Der gemeinsame Fit liefert $\chi^2=4.102$ bei drei Freiheitsgraden; der eingefrorene TFPT-Seed liefert $\chi^2=4.513$ bei vier Freiheitsgraden. Keiner der 14 untersuchten Nachbardecoder ist besser.
\item In der gespeicherten Vergleichsmenge von 2000 zufälligen Decodern ist keiner besser. Das ist ein Ergebnis für diese gewählte explorative Vergleichsmenge, keine Wahrscheinlichkeit dafür, dass TFPT wahr ist.
\item Die vorhandene zurückgehaltene $\beta$-Auslese ergibt $0.241304\pm0.001786$ Grad und liegt mit den gespeicherten Messdaten $0.36$ Standardabweichungen vom Vergleichswert entfernt.
\item Die Baum-/Rückkopplungsvarianten sind noch nicht entschieden: $\chi^2_{\rm tree}=4.213$, $\chi^2_{\rm ret}=4.513$. Der kleine Unterschied von $0.300$ und der gespeicherte Cabibbo-Tail-Test von $0.282$ Standardabweichungen wählen keinen der beiden Ursprünge aus.
\end{itemize}
Diese Zahlen beziehen sich auf die vorhandenen eingefrorenen Daten, nicht auf eine Aktualisierung sämtlicher Messwerte. Der Schluss lautet: Die gemeinsame Seed-/Decoderstruktur ist eine ernsthafte Rückbedingung; ihr Erfolg bestimmt noch nicht sämtliche mikroskopischen Wortprodukte oder den Anfangszustand.

\file{quantum-testbed/README.md} erklärt dagegen ausdrücklich einen internen Konsistenztest: Kernelbelegung und geometrische Ratenleiter gehören zum Aufbau. Das bestätigt bedingte Dynamik, nicht ihre physische Herkunft. \file{qgeo-eit-soff/README.md} unterscheidet reale EIT-Daten zur Instrumentenprüfung von der noch fehlenden entscheidenden Viermarken-Zielmessung. Eine erfolgreiche Kalibrierung und ein korrekt negatives Viermarken-Ergebnis sind keine Bestätigung einer Viermarkenstruktur der Natur.

\subsection*{Der zusammenhängende Weg zur gesuchten Lösung}
Das vorhandene Programm enthält eine passende Vorwärtskette; sie muss nicht durch eine neue freie Matrix ersetzt werden:
\[
\begin{gathered}
 (C_\Sigma,\omega,\tau_t,\Theta,[u_\Sigma])
 \longrightarrow
 \bigl(D_\Sigma(A,g,z),\Omega,h,\nabla_{\det}\bigr)\\
 \longrightarrow
 \text{gemeinsame Eich-, Flavor-, Zeit- und Geometrieantworten}.
\end{gathered}
\]
Die Symbole stehen für ursprünglichen Randkern, Zustand, Zeit, Reflexion und markierte Deckwirkung sowie ihre gemeinsame Operatorfamilie und Determinantenverbindung. Die Variationen nach Eichfeld $A$, Metrik $g$ und Familientransport $z$ dürfen nicht in voneinander unabhängig angepassten Quellen erfolgen. Orientierung und geladene Felder müssen auf demselben Weg erhalten bleiben.

\file{source-three-route-closure-20260922} formuliert bereits eine solche gemeinsame operative Quelle mit markierten Schwingerantworten, Reflexion, Zustand, Zeit und Kragendynamik. \file{SEAM.MMST.TYPEIII.CHARGED.01} nennt den tatsächlichen geladenen Kontinuumsanschluss; \file{SEAM.DETLINE.UNIFICATION.01} verlangt danach den Anschluss an die vierdimensionale Determinantenlinie einschließlich Verbindung und Holonomien. Diese Verträge sind ein zusammenhängender bedingter Weg, keine bereits geschlossene Ursprungsherleitung. Der ausgewählte operative Seed enthält Zeit und Zustand weiterhin als Daten.

\picturetext{Die alten Herleitungen sind gekoppelte Anzeigen derselben gesuchten Maschine. Die neue Prüfung sagt, dass wir bisher zum Teil die Anzeigen für die Maschine gehalten haben. Gesucht wird die ursprüngliche einfache Regel, die auch die verborgenen Übergänge trägt und dadurch alle Anzeigen gemeinsam erzeugt. Eine beliebige zusätzliche Maschine, die jede Anzeige einzeln nachstellt, würde diese Frage nicht lösen.}

Damit wird die Priorität präziser: Zuerst ist der ursprüngliche volle graduierte Wortprozess samt seiner tatsächlichen Zustand-/Zeitregel zu verfolgen. Danach sind \emph{dieselben} geladenen Operatoren unter den schon verlangten Eich-, Familien- und Metrikvariationen zu übertragen. Die bestehenden Alpha-, Massen-, Flavor- und Neutrinoherleitungen sind dabei gleichzeitig einzuhaltende Rückbedingungen. Wo ein gemischtes Produkt, eine Orientierungsantwort oder ein gemeinsamer Anschluss fehlt, liegt eine konkret benannte Herkunftsaufgabe. Die erneute Suche nach einer einzigen neutralen Matrix oder nach einem Spektrum, das alle Sektoren ersetzt, setzt erneut die gerade widerlegte Vereinfachung voraus.

\section*{Wie dieses Dokument zu lesen ist}
\addcontentsline{toc}{section}{Wie dieses Dokument zu lesen ist}
Der Auftrag lautet, die vollständige Verbindung zu finden. Deshalb ist der Maßstab dieses Textes eine gemeinsame Herleitung mit denselben Feldern, demselben Zustand und derselben Dynamik. Eine Sammlung passender Zahlen, ein funktionierendes Hilfsmodell oder eine große Zahl grüner Prüfungen reicht dafür nicht.

Die Lesereihenfolge ist bewusst gestuft: Zuerst kommt das Gesamtbild; anschließend die ursprüngliche Kette und ihre neuen Verbindungen; danach Universalraum, Raum und physische Realisierung; schließlich Anwendungen, Algorithmen und die konkrete Abschlussagenda. Die ausführbare Quellenkonstruktion steht in den Abschnitten \ref{sec:levelselection}--\ref{sec:physicalfollow}. Die anschließende Gesamtprüfung in Abschnitt~\ref{sec:jointclosure} verbindet Ursprung, K3, Quantisierung, lokale Zeit und Teilchenstatistik und korrigiert die bislang zu unspezifische Abschlussagenda. Wer zunächst nur die Bedeutung verstehen möchte, kann die eingerahmten Kernaussagen und Abschnitte \ref{sec:bigpicture}, \ref{sec:jointclosure} und \ref{sec:completion} lesen.

\begin{tabularx}{\linewidth}{lY}
\toprule
Zeichen & Bedeutung in diesem Dokument\\\midrule
\E & Exakte mathematische Aussage oder exakte endliche Rechnung innerhalb ausdrücklich genannter Voraussetzungen. Das Zeichen behauptet keine vollständige Formalisierung in Lean.\\
\N & Numerisch oder durch einen konkreten ausgeführten Test bestätigt; Reichweite durch Eingaben, Genauigkeit und Abdeckung begrenzt.\\
\B & Bedingter Anschluss: Die Folgerung trägt, wenn die benannte Identifikation oder Zusatzannahme erfüllt ist.\\
\Opend & Offene Herleitung, physische Identifikation, Auswahl oder empirische Behauptung.\\
\Xclaim & Eine konkrete Gleichsetzung, Konstruktion oder stärkere Behauptung ist durch einen Gegenfall ausgeschlossen.\\\bottomrule
\end{tabularx}

Die gesamte Sitzung ist hier inhaltlich konsolidiert, nicht als chronologisches Protokoll wiedergegeben. Überholte Zwischenstände bleiben dort sichtbar, wo ihre Korrektur die Interpretation verändert. Die Anhänge benennen die verwendeten Quellenfamilien, den ausgeführten Prüfstand und die weiterhin offenen Anschlüsse. Der dokumentierte Explorer-Schnappschuss vom 28.~September enthält 316 interne Prüfungen; frühere Zwischenstände vom 29.~September enthielten 343 beziehungsweise 367, der Massentransportstand enthielt 372 und die Wand-Prüfung 375. Der Majorana-Operatoranschluss bestand alle 380 internen Prüfungen. Der Stand einschließlich gemeinsamer KMS-/Seesaw-Auswertung bestand alle 381 internen Prüfungen. Der gemeinsame Lauf mit Ladungs-/Flavorwörterbuch und Originaltransport bestand alle 383 internen Prüfungen; mit dem tatsächlichen Chevalley-Abschluss des Doppel-Twists bestand der vollständige Explorer 384 Prüfungen in 24 Stufen. Der gezielte Integrationslauf der Ladungsantwort, des Spinlifts und der Synthese besteht einschließlich Fourierrekursion und Paarkern-Scope 21 Tests. Diese Zahlen bezeichnen weder voneinander unabhängige Beweise noch eine vollständige Prüfung aller Paper.

\tableofcontents
\newpage

\section{Zurück zum Ursprung: Ein allgemeiner Auswahltest für das ganze Gerüst}\label{sec:globaloriginselection}

\key{\textbf{Der neue verbundene Befund lautet:} Die ursprüngliche Zwei-Seiten-Geometrie lässt sich ohne vorab eingesetztes E8 global untersuchen. Widerspruchsfreiheit allein lässt unendlich viele größere Abschlüsse zu. Bereits ein einziger lokaler, spinorgeladener Strom vom konformen Gewicht eins würde ausschließlich die Viermarkenstruktur auswählen; die zuvor verlangte vollständige Stromerzeugung ist dafür nicht nötig. In der vorhandenen Achtkanal-Fermionisierung ist der entsprechende E8-Generator außerdem exakt das führende Feld im renormierten Produkt zweier Vierteltwists. Diese algebraische Verbindung ist konstruiert. Dass der ursprüngliche mikroskopische Nahtprozess genau dieses lokale Feld hervorbringt, bleibt nachzuweisen. Die Auswahl legt vier Marken, Zyklengitterrang drei, Trägerrang fünf und E8 gemeinsam fest; die originale Familien-/Materielesart und ihre physische Realisierung behalten ihre eigenen Voraussetzungen.
}

\subsection{Das einfache Bild und die genaue Frage}

Die ursprüngliche Naht wird auf zwei Arten gelesen. Die eine Seite liest Wege um ihre Marken; die andere liest Funktionen mit erlaubten Polen an denselben Marken. Bei vier Marken ergeben sich drei unabhängige Wegzyklen und fünf Funktionsplätze. Das ist der ursprüngliche Zusammenhang von Familien und Träger, nicht die Einführung zweier unabhängiger Räume.

\picturetext{Man nimmt dasselbe Lochmuster einmal als Wegeplan und einmal als Register von möglichen Anzeigen. Gesucht ist ein Muster, dessen beide Lesarten zu einer vollständigen Mechanik zusammenpassen. Die vier Löcher dürfen dabei nicht schon vor der Prüfung als einzig mögliche Zahl eingesetzt sein.}

Wir prüfen deshalb alle ganzen Markenzahlen $d\ge2$ in der folgenden ausdrücklich benannten Vergleichsklasse. Die Grundfläche ist $\mathbb P^1$ mit einem reduzierten Divisor $\mathcal D$ vom Grad $d$. Die Zyklenseite trägt das ursprüngliche integrale Residuengitter, die Funktionsseite wird durch dieselbe gerade Clifford-Konstruktion gelesen:
\[
\begin{aligned}
 H_1(\mathbb P^1\setminus\mathcal D;\Z)&\cong\Z^{d-1},
 &\text{Zyklengitter}&=A_{d-1},\\
 \dim H^0(\mathbb P^1,\mathcal O(\mathcal D))&=d+1,
 &\text{Trägergitter}&=D_{d+1}.
\end{aligned}
\]
Die Ränge folgen aus Topologie und Riemann--Roch. Die Clifford-Lesart und die normierte integrale Gittermetrik gehören zur Vergleichsklasse; bloße Dimensionsgleichheit würde sie nicht beweisen. Setze
\[
 R_d=A_{d-1}\oplus D_{d+1}.
\]
Gesucht wird ein positives gerades unimodulares Gitter $\mathcal L\supset R_d$. Im Gitter-Randmodell bedeutet dies einen holomorphen Abschluss. Es ist noch keine Aussage über vierdimensionale Raumzeit, physische Teilchenstatistik oder ein Vakuum mit gebrochener Eichsymmetrie.

\subsection{Die gesamte Abschlussklasse ist unendlich -- ein globaler Satz}

\textbf{Satz 1.} Ein solcher Abschluss existiert genau für
\begin{equation}\label{eq:globalgluetower}
 \boxed{d=4m^2,\qquad m=1,2,3,\ldots .}
\end{equation}

\textbf{Notwendigkeit.} Die Determinanten sind $\det A_{d-1}=d$ und $\det D_{d+1}=4$. Für den Erweiterungsindex $r=[\mathcal L:R_d]$ gilt deshalb
\[
 1=\det\mathcal L=\frac{4d}{r^2}.
\]
Also ist $d$ ein Quadrat. Ein positives gerades unimodulares Gitter hat einen durch acht teilbaren Rang. Sein Rang ist hier $2d$, also gilt zusätzlich $4\mid d$. Beide Forderungen ergeben genau $d=4m^2$.

\textbf{Hinreichendheit.} Für $d=4m^2$ ist die Diskriminantengruppe
\[
 A(R_d)=\Z_4\oplus\Z_d.
\]
Bezeichne ihre Generatoren mit $s$ auf der $D$-Seite und $a$ auf der $A$-Seite. In der Normkonvention modulo $2\Z$ gilt
\[
 q_D(s)=\frac{d+1}{4},\qquad q_A(a)=\frac{d-1}{d}.
\]
Die zyklische Untergruppe
\[
 H_m=\langle(s,ma)\rangle
\]
hat Ordnung $4m=\sqrt{4d}$ und
\[
 q(s,ma)=\frac{4m^2+1}{4}
       +m^2\frac{4m^2-1}{4m^2}=2m^2\equiv0\pmod{2\Z}.
\]
Auch jeder Vielfache des Generators ist isotrop. Die zugehörige Überlagerung ist daher gerade und wegen $|H_m|^2=|A(R_d)|$ unimodular. Verwendet wird die klassische Korrespondenz zwischen isotropen Diskriminantenuntergruppen und geraden Überlagerungsgittern~\cite{nikulinforms}; die vorstehende Anwendung auf die ganze $d$-Familie ist hier vollständig ausgeschrieben. $\square$

\begin{center}
\begin{tabular}{rrrrl}
\toprule
$m$ & Marken $d$ & Rang & Index & Zusätzliche Norm-2-Wurzeln\\\midrule
1 & 4 & 8 & 4 & vorhanden; Abschluss E8\\
2 & 16 & 32 & 8 & keine\\
3 & 36 & 72 & 12 & keine\\
4 & 64 & 128 & 16 & keine\\\bottomrule
\end{tabular}
\end{center}
Die Tabelle illustriert den allgemeinen Beweis; sie ersetzt ihn nicht. Ihre Diskriminantenwerte, Isotropie, Indizes und Coset-Minimalnormen wurden zusätzlich in exakter Brucharithmetik ausgeführt. Für $m=2$ ist bereits $D_{17}\oplus A_{15}$ mit $H_2=\langle(s,2a)\rangle$ ein konkreter Gegenfall gegen die Behauptung, gerader unimodularer Abschluss allein erzwinge E8.

Der Originalvertrag \file{E8.ONEOBJECT.01} ist enger: Er verlangt \emph{diagonalen} Glue, also trivialen Schnitt mit beiden einzelnen Diskriminantenfaktoren. Ein solcher voller Glue ist der Graph einer Antiisometrie und verlangt gleiche Gruppenordnungen, daher $d=4$. Sein Census $2\le d\le12$ ist in dieser Klasse korrekt. Die diagonale Bedingung darf aber nicht beim Übergang zu allgemeinem holomorphem Abschluss weggelassen werden. In der obigen größeren Familie hat $H_m$ einen reinen Zyklensektor vom Umfang $m$; erst bei $m=1$ ist dieser trivial.

